\section{Pokyčių valdymas}

\subsection{Tikslas ir apibrėžimas}
Šio skyriaus tikslas – nustatyti formalų ir kontroliuojamą procesą, skirtą bet kokiems projekto apimties, tvarkaraščio, biudžeto ar kokybės reikalavimų pakeitimams valdyti po to, kai pradiniai projekto planai (bazės) yra patvirtinti. 

Pokyčių valdymu siekiama apsaugoti projektą nuo nekontroliuojamo apimties išsiplėtimo (\textit{angl. scope creep}), užtikrinti skaidrų sprendimų priėmimą ir garantuoti, kad visi suinteresuotieji asmenys supranta planuojamų pakeitimų pasekmes.

\subsection{Pokyčių klasifikacija}
Visi projekto pakeitimai skirstomi į tris pagrindines kategorijas pagal jų poveikį:
\begin{itemize}
    \item \textbf{Nedideli pakeitimai:} Įtakos neturi nei projekto biudžetui, nei galutiniam terminui (pvz., nedidelis užduočių eiliškumo pasikeitimas komandos viduje).
    \item \textbf{Vidutiniai pakeitimai:} Keičia tarpinius terminus arba reikalauja išteklių perskirstymo nekeičiant patvirtintos apimties ir bendrų darbo sąnaudų.
    \item \textbf{Dideli pakeitimai:} Keičia projekto apimtį, didina patvirtintas darbo sąnaudas arba lemia galutinio termino pratęsimą.
\end{itemize}

\subsection{Pokyčių valdymo procesas}
Bet koks pokytis projekto eigoje privalo praeiti šiuos 5 etapus:

\begin{enumerate}
    \item \textbf{Pokyčio inicijavimas ir pateikimas} --
    Bet kuris projekto komandos narys ar užsakovas gali pastebėti pokyčio poreikį. Iniciatorius privalo užpildyti formalią \textit{Pokyčio užklausą}, kurioje aprašo pokyčio esmę bei priežastį.
    
    \item \textbf{Registracija} -- 
    Projekto vadovas gauna užklausą, ją užregistruoja centrinėje sistemoje -- \textit{Pokyčių registre} ir suteikia unikalų numerį.
    
    \item \textbf{Poveikio analizė} -- 
    Projekto vadovas kartu su atsakingais specialistais įvertina pokyčio poveikį:
    \begin{itemize}
        \item \textit{Apimčiai:} Kokios naujos funkcijos ar reikalavimai atsiranda?
        \item \textit{Laikui:} Kiek dienų ar savaičių užtruks įgyvendinimas?
        \item \textit{Kaštams:} Kokie bus reikalingi papildomi finansiniai kaštai?
        \item \textit{Rizikai:} Kokias naujas grėsmes ar galimybes tai sukuria?
    \end{itemize}
    
    \item \textbf{Sprendimo priėmimas} -- 
    Remiantis poveikio analize, pokyčio užklausa leidžiama tvirtinti atsakingiems asmenims. Priklausomai nuo pokyčio dydžio, sprendimą priima Projekto vadovas arba Pokyčių valdymo komitetas. Galimi sprendimai: patvirtinti, atmesti arba grąžinti patikslinimui.
    
    \item \textbf{Įgyvendinimas ir atnaujinimas} -- 
    Jei pokytis patvirtinamas, projekto vadovas oficialiai atnaujina projekto planą, biudžetą bei tvarkaraštį ir paveda komandai įgyvendinti užduotį. Atmesti prašymai archyvuojami nurodant atmetimo priežastį.
\end{enumerate}

\subsection{Valdymo kompetencijos ir ribos}
\begin{table}[htbp]
\centering
\caption{Pokyčių tvirtinimo atsakomybės matrica}
\label{tab:pokyciu_matrica}
\begin{tabularx}{\textwidth}{l X l}
\toprule
\textbf{Pokyčio lygis} & \textbf{Kriterijai / Apibūdinimas} & \textbf{Tvirtinantis asmuo} \\ 
\midrule
Žemas & Veiklų perskirstymas be papildomų sąnaudų ir be KT ar galutinio termino pakeitimo & Projekto vadovas \\
Vidutinis & Keičiamas tarpinis KT, bet ne apimtis, bendros sąnaudos ar galutinis terminas & Pokyčių valdymo komitetas \\
Aukštas & Keičiama apimtis, didinamos bendros sąnaudos arba keičiamas galutinis terminas & Užsakovas, gavęs komiteto išvadą \\
\bottomrule
\end{tabularx}
\end{table}

Rizikų rezervą (26~žm.~d.) projekto vadovas naudoja tik registre
įvardytoms rizikoms, fiksuodamas panaudojimą pagal rizikos ID.
Valdymo rezervo (37~žm.~d.) panaudojimą tvirtina užsakovas. Naujos
funkcijos ir išplėtimai vertinami bei tvirtinami atskirai; jų sąnaudos
automatiškai nedengiamos iš rezervų.
